% Vietnamese translation for OI Public Library
% Source-File: IOI中国国家候选队论文/2008/Day1/8.吕子鉷《浅谈最短径路问题中的分层思想》/浅谈最短路径问题中的分层思想.doc
\documentclass[11pt]{article}
\usepackage{oipl-vn}

\newcommand{\Oh}{\mathcal{O}}

\title{Bàn về tư tưởng phân tầng trong bài toán đường đi ngắn nhất}
\author{Lu Zihong\\Trường Trung học số 7 Tuyền Châu, Phúc Kiến}
\date{Trại đông Tin học toàn quốc, 2008}

\begin{document}
\maketitle

\origtitle{Bàn về tư tưởng phân tầng trong bài toán đường đi ngắn nhất}
\sourcepath{IOI China candidate team papers / 2008 / Day 1 / Lu Zihong / layered shortest paths.doc}

\renewcommand{\abstractname}{Tóm tắt}
\begin{abstract}
Tư tưởng phân tầng và thuật toán đường đi ngắn nhất đều có phạm vi ứng dụng
rộng. Bài viết phân tích một số đề thi tin học để trình bày hai cách dùng
phân tầng trong bài toán đường đi ngắn nhất: dựng mô hình sao cho bài toán
phức tạp trở thành một bài toán đường đi ngắn nhất thông thường, và khai thác
cấu trúc tầng để tối ưu thuật toán. Các ví dụ gồm \textit{Rescue}, \textit{Fence},
\textit{Cow Relay}, \textit{Bicriterial Routing} và \textit{Roads}.
\end{abstract}

\noindent\textbf{Từ khóa:} đường đi ngắn nhất; phân tầng; lý thuyết đồ thị.

\section{Dẫn nhập}

Bài toán đường đi ngắn nhất là một bài toán kinh điển của lý thuyết đồ thị.
Trọng số cạnh không nhất thiết chỉ là khoảng cách hình học; trong nhiều mô
hình, nó có thể là thời gian, chi phí hoặc một đại lượng cộng dồn tuyến tính
dọc theo đường đi. Vì vậy đường đi ngắn nhất xuất hiện tự nhiên trong quy
hoạch đô thị, định tuyến giao thông, tối ưu mạng và nhiều bài thi tin học.

Tư tưởng phân tầng cũng là một công cụ quan trọng. Ta gặp nó trong cách chia
giai đoạn của quy hoạch động, trong các thuật toán luồng cực đại dựa trên đồ
thị mức, và trong nhiều mô hình trạng thái khác. Khi kết hợp phân tầng với
đường đi ngắn nhất, ta có thể đưa những ràng buộc tưởng như nằm ngoài đồ thị
vào cấu trúc đỉnh và cạnh của một đồ thị mới.

Bài viết xét hai hướng chính. Phần đầu dùng phân tầng để dựng mô hình. Phần
sau dùng tính chất của đồ thị phân tầng để tối ưu thuật toán.

\section{Dùng phân tầng để dựng mô hình}

Trong phần này, tác giả xét ba ví dụ: \textit{Rescue}, \textit{Fence} và
\textit{Cow Relay}. Điểm chung là bài toán ban đầu không phải một bài toán
đường đi ngắn nhất chuẩn trên một đồ thị cố định, nhưng có thể trở thành như
vậy nếu mỗi tầng biểu diễn một điều kiện bổ sung.

\subsection{Ví dụ 1: Rescue}

Một mê cung hình chữ nhật gồm \(N\) hàng và \(M\) cột. Hai ô kề nhau theo
hướng bắc-nam hoặc đông-tây có thể thông nhau, có thể bị ngăn bởi một bức
tường, hoặc có thể có một cánh cửa khóa. Các cửa được chia thành \(P\) loại;
cùng một loại dùng cùng một chìa khóa, các loại khác nhau dùng chìa khác
nhau. Một số ô chứa chìa khóa. Mỗi bước đi sang ô kề tốn thời gian \(1\);
thời gian nhặt chìa và mở cửa được bỏ qua. Cần tìm thời gian ngắn nhất từ
\((1,1)\) tới \((N,M)\). Giới hạn là \(N,M\le 15\), \(P\le 10\).

Nếu bỏ qua cửa và chìa khóa, mỗi ô là một đỉnh, mỗi cặp ô thông nhau nối bởi
một cạnh trọng số \(1\), và bài toán trở thành đường đi ngắn nhất chuẩn. Cửa
và chìa khóa làm cho khả năng đi qua cạnh phụ thuộc vào trạng thái đã nhặt
chìa. Vì vậy ta chia đồ thị thành \(2^P\) tầng, mỗi tầng ứng với một tập chìa
khóa đang có.

Trong một tầng, các cạnh đi qua cửa chỉ tồn tại nếu tập chìa của tầng đó cho
phép mở cửa. Với một ô chứa chìa, từ đỉnh tương ứng trong tầng hiện tại ta nối
một cạnh trọng số \(0\) tới cùng ô ở tầng sau khi đã thêm chìa khóa ấy. Trên
đồ thị mới, mọi trạng thái đã được biểu diễn bằng đỉnh, nên có thể chạy trực
tiếp thuật toán đường đi ngắn nhất.

Vì mọi cạnh có trọng số \(0\) hoặc \(1\), ta có thể dùng BFS hai đầu hàng đợi
hoặc một biến thể BFS cho trọng số \(0/1\). Độ phức tạp thời gian và bộ nhớ
đều là
\[
  \Oh(2^P NM).
\]

\subsection{Ví dụ 2: Fence}

Một điền trang có tòa nhà chính và nhiều công trình khác. Chủ điền trang muốn
xây một hàng rào bao quanh tòa nhà chính, nhưng hàng rào không được đi quá
gần bất kỳ công trình nào. Mỗi công trình được bao trong một hình chữ nhật
cấm, các cạnh song song trục tọa độ; hàng rào cũng chỉ được gồm các đoạn song
song trục. Cần tìm độ dài nhỏ nhất của một hàng rào hợp lệ bao quanh tòa nhà
chính. Số hình chữ nhật \(m\le 100\), tọa độ không âm và không vượt quá
\(10000\).

Mục tiêu là tối thiểu hóa một đại lượng cộng dồn tuyến tính, nên ta thử mô
hình hóa bằng đường đi ngắn nhất. Trước hết xây một đồ thị một tầng: các đỉnh
là các đỉnh của các hình chữ nhật. Nếu đoạn thẳng song song trục nối hai đỉnh
không đi vào bên trong hình chữ nhật cấm nào, ta nối hai đỉnh ấy bằng một cạnh
có trọng số bằng khoảng cách Manhattan. Mỗi chu trình trong đồ thị tương ứng
với một phương án hàng rào, và độ dài chu trình là độ dài hàng rào.

Tuy nhiên, không phải chu trình nào cũng hợp lệ; nó phải bao quanh tòa nhà
chính. Theo tiêu chuẩn hình học quen thuộc, một điểm nằm trong đa giác nếu một
tia xuất phát từ điểm ấy cắt biên đa giác số lẻ lần. Vì vậy ta kẻ một tia từ
tòa nhà chính và chỉ cần biết chu trình cắt tia ấy với số lần chẵn hay lẻ.

Ta tách đồ thị thành hai tầng, biểu diễn tính chẵn lẻ của số lần cắt tia. Nếu
một cạnh của đồ thị một tầng cắt tia, cạnh đó nối hai tầng khác nhau; nếu
không, nó nối trong cùng tầng. Khi đó tìm một chu trình bao quanh tòa nhà
chính tương đương với tìm đường đi ngắn nhất giữa hai bản sao khác tầng của
cùng một đỉnh. Lấy nhỏ nhất trên mọi đỉnh sẽ cho độ dài hàng rào tối ưu.

Nếu kiểm tra mọi cặp đỉnh bằng liệt kê hình chữ nhật, việc dựng cạnh tốn
\(\Oh(m^3)\). Chạy Dijkstra từ các đỉnh cần xét cũng có thể làm trong
\(\Oh(m^3)\). Bộ nhớ là \(\Oh(m^2)\).

\subsection{Ví dụ 3: Cow Relay}

Farmer John chọn \(k\) con bò tham gia tiếp sức. Trang trại có \(n\) cánh
đồng và \(m\) cạnh vô hướng có trọng số, biểu diễn thời gian đi qua. Các con
bò lần lượt chạy từ cánh đồng \(1\) tới cánh đồng \(n\); khi một con tới đích,
con sau lập tức xuất phát. Đường đi của hai con bất kỳ phải khác nhau. Cần tìm
thời điểm con bò thứ \(k\) tới đích. Giới hạn là \(k\le 40\), \(n<800\),
\(m\le 4000\).

Nhìn theo đồ thị, bài toán yêu cầu tìm đường đi thứ \(k\) ngắn nhất giữa hai
đỉnh trong đồ thị có trọng số dương. Một cách trực tiếp là duy trì hàng đợi ưu
tiên của các đường đi. Mỗi lần lấy đường đi ngắn nhất hiện có, mở rộng nó theo
các cạnh ra từ đỉnh cuối, cho tới khi thu được \(k\) đường từ \(1\) tới \(n\).
Vì mỗi đỉnh chỉ cần nhiều nhất \(k\) đường đi ứng viên để có thể thuộc một
trong \(k\) đường tốt nhất tới \(n\), số đường được mở rộng không vượt quá
\(kn\). Độ phức tạp có thể viết là
\[
  \Oh(km+kn\log(kn)).
\]
Cách này đúng nhưng chưa đủ đẹp.

Bản gốc đưa ra một cách nhìn khác. Gọi \(P\) là tập mọi đường đi từ \(1\) tới
\(n\), và \(p_i\) là đường đi ngắn thứ \(i\). Nếu sau khi tìm được \(p_i\) ta
loại bỏ đường đó khỏi tập đường đi, thì đường ngắn nhất còn lại chính là
\(p_{i+1}\). Vấn đề là phải thực hiện thao tác ``loại một đường đi'' mà không
phá hỏng toàn bộ đồ thị.

Ta thêm hai đỉnh \(s,t\), cạnh \((s,1)\) và \((n,t)\) có trọng số \(0\). Mỗi
đường từ \(1\) tới \(n\) tương ứng một đường từ \(s\) tới \(t\). Sau khi tìm
đường ngắn nhất \(p_{st1}\), sao chép một phần các đỉnh và cạnh của đồ thị.
Với mỗi cạnh \((i,j)\) không thuộc \(p_{st1}\), ta thêm cạnh chuyển từ bản gốc
của \(i\) sang bản sao của \(j\). Với các cạnh thuộc \(p_{st1}\), không thêm
cạnh chuyển như vậy. Các cạnh trong tầng sao chép giữ nguyên trọng số.

Nếu xem mỗi đỉnh gốc và bản sao của nó là cùng một đỉnh, mọi đường từ \(s\)
tới \(t'\) trong đồ thị mới đều là một đường cũ từ \(s\) tới \(t\), nhưng
không thể là đúng \(p_{st1}\), vì đường ấy bị chặn tại cạnh đầu tiên cần
chuyển tầng. Ngược lại, mọi đường khác \(p_{st1}\) đều có cạnh đầu tiên khác
với \(p_{st1}\), nên tại cạnh đó có thể chuyển sang tầng sao chép và tiếp tục
tới \(t'\). Vì vậy tập đường từ \(s\) tới \(t'\) chính là tập cũ sau khi loại
\(p_{st1}\). Đường ngắn nhất trong tập này là đường ngắn thứ hai.

Lặp quá trình trên sẽ thu được các đường ngắn tiếp theo. Không cần sao chép
toàn bộ đồ thị ở mỗi bước: nếu \(v_k\) là đỉnh đầu tiên trên \(p_{st1}\) có
bậc vào lớn hơn \(1\), chỉ cần sao chép đoạn từ \(v_k\) trở đi. Những đỉnh
không nằm trên đường hoặc nằm trước \(v_k\) không tạo ra đường chuyển tầng hữu
ích từ \(s\).

Do đồ thị mới chỉ thêm đỉnh và cạnh vào đồ thị cũ, cây đường đi ngắn nhất của
phần cũ không bị ảnh hưởng. Khi tìm đường đi mới, ta có thể chỉ duy trì cây
đường đi ngắn nhất trên các đỉnh vừa sao chép, thay vì chạy lại toàn bộ. Nếu
dùng SPFA cho cây đường đi ngắn nhất, mỗi vòng có thể xem là \(\Oh(m)\), tổng
thời gian \(\Oh(km)\), bộ nhớ \(\Oh(kn+m)\).

\subsection{Nhận xét về dựng mô hình}

Ba ví dụ trên dùng phân tầng theo ba cách khác nhau. Trong \textit{Rescue},
tầng biểu diễn trạng thái chìa khóa, tức điều kiện quyết định một số cạnh có
đi được hay không. Trong \textit{Fence}, tầng biểu diễn tính chẵn lẻ của số
lần cắt một tia, tức điều kiện để chu trình bao quanh tòa nhà chính. Trong
\textit{Cow Relay}, tầng và cạnh chuyển tầng dùng để loại bỏ một đường đi đã
chọn, làm cho mỗi tầng có một tập đường đi khác nhau.

Bản chất là: khi một đồ thị một tầng không biểu diễn đủ các điều kiện của bài
toán, ta sao chép nó thành nhiều tầng, mỗi tầng tương ứng một điều kiện. Việc
phân tầng làm tăng kích thước đồ thị, nhưng nếu khai thác được cấu trúc của
các tầng thì phần tăng thêm thường vẫn chấp nhận được.

\section{Dùng phân tầng để tối ưu thuật toán}

Phần trước dùng phân tầng để biến bài toán thành đường đi ngắn nhất. Phần này
xét chiều ngược lại: khi đã có đồ thị phân tầng, chính cấu trúc tầng có thể
giúp giảm độ phức tạp.

\subsection{Ví dụ 4: Bicriterial Routing}

Mạng đường thu phí của Byteland có các đường hai chiều. Mỗi đường có thời gian
đi và phí phải trả. Với cùng điểm xuất phát và đích, một tuyến đường \(A\) tốt
hơn tuyến \(B\) nếu \(A\) nhanh hơn mà phí không cao hơn, hoặc phí thấp hơn mà
không chậm hơn. Một tuyến được gọi là tối ưu theo hai tiêu chí nếu không có
tuyến nào tốt hơn nó. Cần đếm số cặp \((\text{phí},\text{thời gian})\) tối ưu
từ \(s\) tới \(e\). Giới hạn: \(n\le 100\), \(m\le 300\), mỗi phí và thời gian
không vượt quá \(100\).

Khác với đường đi ngắn nhất thường, mỗi cạnh có hai trọng số. Ta chọn phí làm
chỉ số tầng. Nếu \(c\) là phí lớn nhất của một cạnh, một đường tối ưu không
cần có tổng phí vượt quá \(nc\) trong mô hình của bài. Do đó ta tách đồ thị
thành \(nc+1\) tầng, tầng \(q\) biểu diễn đã dùng tổng phí \(q\). Một cạnh có
phí \(w\) và thời gian \(t\) nối từ tầng \(q\) sang tầng \(q+w\), với trọng số
đường đi là \(t\).

Đồ thị mới có \(\Oh(nc)\) tầng, mỗi tầng có \(n\) đỉnh và \(m\) cạnh kiểu
tương ứng, nên số đỉnh là \(\Oh(n^2c)\), số cạnh là \(\Oh(nmc)\). Vì thời gian
không âm, có thể chạy Dijkstra trên toàn đồ thị. Cách trực tiếp tốn khoảng
\[
  \Oh(ncm\log(n^2c)).
\]

Nhưng đồ thị có tính chất quan trọng: phí không âm, nên cạnh chỉ đi trong cùng
tầng hoặc từ tầng thấp sang tầng cao hơn. Tầng phí cao không ảnh hưởng ngược
lại tầng phí thấp. Vì vậy thay vì chạy một Dijkstra trên toàn bộ đồ thị, ta
xử lý từng tầng theo thứ tự phí tăng dần; mỗi tầng chỉ cần một lần đường đi
ngắn nhất trên \(n\) đỉnh. Nếu dùng heap nhị phân, mỗi tầng tốn
\(\Oh(m\log n)\), tổng cộng
\[
  \Oh(ncm\log n).
\]

Sau khi có thời gian tốt nhất cho từng mức phí, ta quét các cặp
\((\text{phí},\text{thời gian})\) để loại những cặp bị trội và đếm các cặp
tối ưu.

\subsection{Ví dụ 5: Roads}

Có \(N\) thành phố và các đường một chiều. Mỗi đường có chiều dài và phí. Bob
có nhiều nhất \(k\) đồng, xuất phát từ thành phố \(1\), muốn tới thành phố
\(N\) nhanh nhất trong phạm vi chi phí cho phép. Giới hạn: \(k\le 10000\),
\(n\le 100\), \(m\le 10000\), chiều dài mỗi đường không vượt quá \(100\), phí
mỗi đường không âm và không vượt quá \(100\).

Mô hình đầu tiên giống bài trước: tách đồ thị thành \(k+1\) tầng theo số tiền
đã dùng. Một cạnh có phí \(w\) nối từ tầng \(q\) sang tầng \(q+w\), trọng số
là chiều dài đường. Vì chiều dài dương, có thể dùng Dijkstra. Nếu cài đặt đơn
giản bằng mảng cho đồ thị khá dày, độ phức tạp là
\[
  \Oh(k(kn^2+m)).
\]

Tương tự \textit{Bicriterial Routing}, vì phí không âm nên có thể xử lý các
tầng theo thứ tự phí tăng dần, giảm còn
\[
  \Oh(k(n^2+m)).
\]

Tuy vậy, bài này có thêm điều kiện mà mô hình theo phí chưa dùng hết: chiều
dài đường là số nguyên dương và có chặn nhỏ. Ta có thể đổi trục phân tầng sang
chiều dài và dùng quy hoạch động. Gọi \(f[i,j]\) là số tiền ít nhất cần dùng
để tới thành phố \(j\) bằng một đường có tổng chiều dài đúng \(i\). Khi có
cạnh từ \(j_0\) tới \(j\) với chiều dài \(\ell\) và phí \(w\), ta cập nhật
\[
  f[i,j] =
  \min\bigl(f[i,j],\, f[i-\ell,j_0]+w\bigr).
\]
Điều kiện đầu là \(f[0,1]=0\), các trạng thái khác ban đầu là vô hạn. Đáp án
là \(i\) nhỏ nhất sao cho \(f[i,N]\le k\). Nếu \(L\) là chiều dài cạnh lớn
nhất, cách này có độ phức tạp \(\Oh(nLm)\), tốt hơn trong bối cảnh giới hạn
của đề.

\subsection{Nhận xét về tối ưu}

Đồ thị phân tầng thường có nhiều phần giống nhau giữa các tầng. Ta không nhất
thiết phải lưu toàn bộ đỉnh và cạnh một cách tường minh; nhiều cạnh có thể
được sinh khi cần. Ngoài ra, nếu cạnh chỉ đi từ tầng thấp sang tầng cao, các
tầng tạo thành một thứ tự topo theo chiều tầng. Khi đó ta có thể xử lý từng
tầng, làm giảm kích thước hàng đợi ưu tiên hoặc số vòng lặp của Bellman-Ford
và SPFA.

Trong trường hợp cấu trúc trong cùng một tầng cũng đơn giản, đường đi ngắn
nhất có thể được thay bằng quy hoạch động. So với Dijkstra, quy hoạch động bỏ
được thao tác trên hàng đợi ưu tiên, nên hằng số nhỏ hơn và cài đặt rõ hơn.

\section{Kết luận}

Bài viết trình bày hai vai trò của phân tầng trong bài toán đường đi ngắn
nhất. Khi dựng mô hình, phân tầng giúp biểu diễn các yếu tố khó đặt vào một
đồ thị đơn: trạng thái chìa khóa, tính chẵn lẻ của giao điểm, hay tập đường
đi còn lại sau khi loại một đường. Khi tối ưu thuật toán, cấu trúc tầng giúp
xử lý theo thứ tự tự nhiên, tiết kiệm không gian và thời gian, thậm chí thay
đường đi ngắn nhất bằng quy hoạch động.

Phân tầng không phải một công thức máy móc. Cần chọn đại lượng làm tầng sao
cho nó thật sự nắm giữ yếu tố gây khó của bài toán; đồng thời phải khai thác
được tính chất của tầng để phần phình ra của đồ thị không làm thuật toán mất
hiệu quả.

Tác giả cảm ơn cô Xie Shuiying đã hướng dẫn và hỗ trợ, đồng thời cảm ơn huấn
luyện viên Liu Rujia đã đưa ra nhiều ý kiến quý báu.

\section*{Tài liệu tham khảo}

\begin{enumerate}
  \item Yu Yuanming, \emph{Thuật toán đường đi ngắn nhất và ứng dụng}, bài
  viết đội tuyển quốc gia Trung Quốc năm 2006.
  \item Xiao Tian, \emph{Tư tưởng đồ thị phân tầng và ứng dụng trong thi tin
  học}, bài viết đội tuyển quốc gia Trung Quốc năm 2004.
  \item Baltic Olympiad in Informatics 2002/2007 official solutions.
  \item E. Q. V. Martins và J. L. E. Santos, ``A new shortest paths ranking
  algorithm'', \emph{Investigacao Operacional}, 20(1):47--62, 2000.
\end{enumerate}

\section*{Ghi chú về phụ lục nguồn}

Tài liệu gốc kèm các đề bài đầy đủ của \textit{Rescue}, \textit{Fence},
\textit{Cow Relay}, \textit{Bicriterial Routing} và \textit{Roads}. Phần phụ
lục trong DOC có nhiều đoạn tiếng Anh bị hỏng mã hóa; bản dịch này giữ nội
dung cần cho mạch phân tích trong các ví dụ ở trên và mô tả phụ lục bằng lời
thay vì sao chép phần hỏng.

\end{document}
